\section{서론}\label{sec:introduction}
팔레트의 위치와 방향을 인지하는 능력은 제조·물류 현장에서 화물에 접근하고 포크를 정렬하여 상·하차를 수행하는 데 필요하다. 영상 기반 포크 안내\cite{seelinger2006}, RGB-D 카메라를 이용한 인식·위치 추정\cite{xiao2017}, 반자동 픽업 지원\cite{molter2019} 등은 이 인지 정보를 화물 취급에 연결해 왔다. 자율 화물 취급의 다양한 센서·제어 구성을 정리한 연구\cite{kohne2026}에서도 인지와 작업 실행의 연계가 중요한 문제로 다뤄진다.

단안 컬러 영상으로 팔레트를 검출하는 것과 정확한 자세를 얻는 것은 같지 않다. 외형·시점·조명·가림 변화는 코너의 위치와 대응을 불확실하게 만들고, 검출이 성공한 경우에도 남은 위치 오차가 최종 자세의 정확도를 제한할 수 있다. 팔레트 앞면의 표현을 이용하는 접근\cite{kai2025}, 가림에 대응하는 접근\cite{vu2024}, 합성 자료에서 기하 단서를 학습하는 접근\cite{beleznai2025}은 서로 다른 관측·표현 조건에서 이 문제를 다룬다. 본 연구의 대상은 물체 검출 자체의 실패보다, 이미 얻은 코너를 추가 정보로 수정할 수 있는 범위이다.

인지 정확도를 높이는 방법에는 영상 적응, 초기 추정기의 추가 학습, 추정 이후의 보정이 있다. 이들은 사용하는 감독과 바꾸는 계산 단계가 다르다. 영상 적응과 다단계 자세 보정을 결합한 팔레트 연구\cite{elmoaqet2025}가 존재하므로 후처리의 존재 자체를 새로움으로 삼을 수는 없다. 본 연구는 기존 추정기를 유지하면서 그 내부 영상 특징과 알려진 치수를 함께 사용해 코너를 국소적으로 보정하는 경우를 분리하여 평가한다. 초기 추정기 업데이트는 대표적인 대안으로 비교하되, 자기학습의 모든 성립 조건을 밝히는 것을 연구 범위로 삼지 않는다.

제안 보정기는 초기 코너 주변의 다중 해상도 특징을 읽고, 물리적 치수에서 얻은 정보로 후보 이동의 점수를 조정한다. 점수 분포의 기댓값으로 좌표를 수정하고 원본 영상 좌표에서 이동량을 제한한다. 코너 이외의 객체 선택·박스·검출 점수·중심점은 유지하며, 보정 전후의 자세 계산에 같은 치수와 카메라 정보를 사용한다. 치수 입력은 코너가 주어진 직육면체의 투영을 반드시 만족하도록 강제하는 제약이 아니라, 후보를 평가하는 학습 신호이다.

검증은 세 질문에 대응한다. 첫째, 기존 추정기에 보정을 더하면 코너와 최종 위치·회전이 함께 개선되는가. 둘째, 영상 기반 국소 보정에 치수와 회전 대응 감독을 추가한 효과가 각각 존재하는가. 셋째, 동일 보정 절차를 서로 다른 추정기에 학습·적용했을 때 효과와 실패 양상이 어떻게 달라지는가. 이를 위해 세 기반 추정기의 보정 전후 비교, 구성 요소 절제, 다른 보정 방법과의 비교, 형상·재질·가림별 분석을 수행한다. 결과는 개선의 중심부뿐 아니라 큰 오류, 손상, 실패율과 추가 계산 비용을 함께 보고한다.
